\documentclass[10pt]{article}
\usepackage{icmlko}

\icmlmeta{IP 파운데이션 월드모델}{303}{빌리는 것은 둘뿐이다}{도메인마다 빌린 양을 쟀다 --- 도서가 제 점수의 35\%, 펀딩이 44\%를 빌리고 나머지 여섯은 0 이다}

\begin{document}

\twocolumn[
\icmlheader

\begin{abstract}
노트 301이 \textbf{전이는 자료가 아니라 축을 통해 일어난다}는 것을 보이고
``도메인마다 얼마를 빌리나''를 남겼다. 쟀다 --- 각 도메인의 ``안 오른다''
축(자기 라벨과 주변 관계가 없는 축)을 \textbf{그 도메인에서만} 끄고 그
도메인의 유보 $\rho$ 가 얼마나 내려가는지 본다. \textbf{빌리는 것은 둘뿐이다.}
\textbf{도서 $-$0.1381}(제 점수의 \textbf{35\%}, $t{=}{-}2.87$) ·
\textbf{펀딩 $-$0.1069}(\textbf{44\%}, $t{=}{-}1.97$). 나머지 여섯은
$-$0.0004 $\sim$ $+$0.0098 로 \textbf{0 이거나 오히려 조금 오른다}. 그리고
그 둘은 노트 300이 ``새로운 축이 0개''로 지목한 바로 그 도메인이다.
\textbf{그러니 전이는 보편적인 것이 아니라 대타다} --- 제 신호가 있는
도메인은 안 빌리고, 없는 도메인은 점수의 3분의 1에서 절반을 빌린다.
크기로 보면 이 판의 파운데이션 효과는 \textbf{유보 2,669행 중 243행(9\%)}
위에서만 일어난다. 그래서 노트 301의 ``첫 직접 증거''는 여전히 옳지만
\textbf{좁다} --- 판 무게의 4분의 3을 지는 웹툰 · 애니 · 게임은 아무것도
안 빌린다. 덤으로 아이돌은 끄면 $\rho$ 가 0.0751 에서 0.1632 로
\emph{두 배}가 되는데($t{=}1.53$) 유보 25행이라 못 가른다.
\end{abstract}
]

\section{빌린 양을 정의한다}

노트 301이 ``안 오른다'' 47자리를 통째로 끄면 판이 $-$0.0121 내려가고
그것이 \textbf{국지적}이라는 것을 보였다(도서만 끄면 도서만 무너진다).
그러면 도메인마다 따로 끄고 재면 \textbf{빌린 양}이 나온다.

\begin{center}\small
빌린 양 $\;=\;\rho_d(\text{원래}) - \rho_d(\text{자기 ``안 오른다'' 축 끔})$
\end{center}

값은 0.5 로, 표시자는 0 으로 만든다 --- ``이 도메인은 그 축을 모른다''다.
판정 근거가 학습 구간이므로 유보를 안 보고 고를 수 있다(노트 239 · 242).

\section{둘뿐이다}

\begin{figure*}[t]
\centering
\includegraphics[width=\textwidth]{figs/borrow.pdf}
\caption{왼쪽 --- 자기 ``안 오른다'' 축을 껐을 때 도메인 $\rho$ 의 변화
($\pm$2$\times$짝SE, 부트스트랩 400). 오른쪽 --- 학습행 대 빌린 양.}
\end{figure*}

\begin{center}\scriptsize
\begin{tabular}{lrrrrrr}
\hline
도메인 & 유보 & 축 & 원래 & 막음 & 빌림 & 몫 \\
\hline
\textbf{도서} & 163 & 9 & 0.3988 & 0.2606 & \textbf{0.1381} & \textbf{35\%} \\
\textbf{펀딩} & 80 & 9 & 0.2444 & 0.1375 & \textbf{0.1069} & \textbf{44\%} \\
\hline
애니 & 606 & 2 & 0.5012 & 0.4975 & 0.0037 & 1\% \\
세계애니 & 300 & 4 & 0.5336 & 0.5340 & $-$0.0004 & 0\% \\
모바일 & 441 & 5 & 0.5435 & 0.5505 & $-$0.0070 & $-$1\% \\
웹툰 & 711 & 4 & 0.4619 & 0.4706 & $-$0.0087 & $-$2\% \\
게임 & 180 & 6 & 0.6214 & 0.6312 & $-$0.0098 & $-$2\% \\
아이돌 & 25 & 4 & 0.0751 & 0.1632 & $-$0.0881 & $-$117\% \\
\hline
\end{tabular}
\end{center}

$t$ 는 도서 $-$2.87 · 펀딩 $-$1.97 · 나머지 전부 $|t|{<}1.6$ 이다.
도메인 문턱($|t|{>}2.77$)을 넘는 것은 \textbf{도서 하나}이지만 크기로는
\textbf{둘이 뚜렷하다.}

\paragraph{빌리는 둘이 누구인가} 노트 300이 채택 검사 ③ 을 판 전체에 돌려
\textbf{도서 · 펀딩 · 시장팝업 · 팝업}에 ``새로운 축이 0개''라고 적었다.
지금 빌리는 것이 \textbf{그 도서와 펀딩}이다(시장팝업 · 팝업은 축이 전부
``못 잰다''라 이 실험에 안 들어온다).

\begin{center}\small
\begin{tabular}{ll}
\hline
제 신호가 있는 도메인 & 안 빌린다 (0\% $\sim$ $-$2\%) \\
제 신호가 없는 도메인 & \textbf{35\% $\sim$ 44\% 빌린다} \\
\hline
\end{tabular}
\end{center}

\textbf{전이는 보편적인 것이 아니라 대타다.}

\section{그러면 이 판의 파운데이션 효과는 얼마나 큰가}

\begin{center}\footnotesize
\begin{tabular}{lr}
\hline
빌리는 도메인의 유보 행 & 243 \\
판 전체 채점 행 & 2,669 \\
\hline
\textbf{몫} & \textbf{9\%} \\
\hline
판 무게 상위 셋(웹툰·애니·모바일) & 66\% \\
그 셋이 빌리는 양 & \textbf{0} \\
\hline
\end{tabular}
\end{center}

노트 301의 ``파운데이션 주장의 첫 직접 증거''는 여전히 옳다 --- 다만
\textbf{좁다.} 이 판에서 도메인 가로지르기가 실제로 일하는 자리는 유보의
9\%이고, 무게를 지는 큰 도메인들은 \textbf{자기 자료로만 돈다.}

\paragraph{그리고 이것이 크로스 도메인 갈래를 다시 연다} 노트 291 · 293 ·
294가 크로스 도메인 \emph{자료}(같은 IP · 같은 별칭 · 같은 회사)를 찾다가
전부 문턱 아래로 끝났다. 그 셋은 전부 \textbf{게임 · 모바일 · 만화 · 웹툰}
같은 큰 도메인을 대상으로 했다. \textbf{빌림이 일어나는 자리는 거기가
아니었다} --- 도서와 펀딩이다.

\paragraph{덤 --- 아이돌은 끄면 두 배가 된다} $\rho$ 0.0751 $\to$ 0.1632.
다만 유보 25행에 짝SE 0.0576 이라 $t{=}1.53$, \textbf{못 가른다.}
아이돌은 노트 274가 잰 검출 문턱이 판에서 제일 높은 도메인이다.

\paragraph{남은 것}
\begin{center}\small
\begin{tabular}{ll}
\hline
갈래 & 왜 \\
\hline
도서 · 펀딩에 축을 더 & 빌릴 자리가 있다는 뜻 \\
누구에게서 빌리나 & 도메인을 하나씩 빼서 \\
아이돌 유보를 늘려라 & 두 배가 진짜인지 \\
큰 도메인은 왜 안 빌리나 & 제 신호가 있어서인가 \\
\hline
\end{tabular}
\end{center}

\paragraph{규약} \textbf{전이가 있다고 말할 때 어디서 있는지 같이 적어라.}
노트 301이 판 수준에서 $-$0.0121 을 보고 ``전이는 축을 통해 일어난다''로
끝냈는데, 도메인으로 갈라 보니 \textbf{둘뿐}이고 그 둘이 판의 9\%다.
\textbf{합친 수는 어디서 왔는지를 안 말한다} --- 오늘 노트 298에서 하네스
셋을 합쳤을 때와 같은 함정이고, 그때는 부호가 뒤집힌 셋이 섞여 있었다.

\end{document}
